\chapter{Rigueur Git}
\source{scripts/check-branch.sh, check-commit-msg.sh, check-commit-allowed.sh, install-hooks.sh}
\section{Les règles}
\begin{itemize}
  \item \textbf{Une branche par fonctionnalité ou correction} : \cmd{feat/SPEC-NNN-slug} (par exemple \cmd{feat/SPEC-003-login}) ou \cmd{fix/slug}. Jamais de code sur \cmd{main}.
  \item \textbf{Commits petits et atomiques}, un sujet par commit, messages en anglais au format Conventional Commits : \cmd{feat}, \cmd{fix}, \cmd{refactor}, \cmd{docs}, \cmd{test}, \cmd{chore}.
  \item Le corps du commit cite \cmd{SPEC-NNN}.
  \item Aucun secret, aucun fichier généré lourd, aucun \cmd{node_modules} dans les commits.
\end{itemize}

\section{Les scripts}
\noindent\begin{tabularx}{\linewidth}{>{\raggedright\arraybackslash}p{5.6cm}X}
\toprule
\textbf{Script} & \textbf{Rôle} \\ \midrule
\cmd{check-branch.sh} & refuse \cmd{main}, exige \cmd{feat/...} ou \cmd{fix/...} \\
\cmd{check-commit-msg.sh} & Conventional Commits, 72 caractères maximum sur la première ligne \\
\cmd{check-commit-allowed.sh} & sur \cmd{feat/SPEC-NNN-slug}, refuse le code tant que la spec n'est pas \emph{Validée} \\
\cmd{install-hooks.sh} & installe les hooks \cmd{pre-commit} et \cmd{commit-msg} \\
\bottomrule
\end{tabularx}

\section{Le hook pre-commit}
Installation, sur confirmation (le script modifie le dépôt) : \cmd{scripts/install-hooks.sh}. Il n'écrase jamais un hook existant.
\begin{itemize}
  \item Un commit qui ne touche que \cmd{docs/}, \cmd{.github/}, \cmd{scripts/}, \cmd{tests/} ou des fichiers \cmd{.md} passe toujours.
  \item Les branches \cmd{fix/...} ne sont pas soumises à la règle de spec.
  \item Une branche \cmd{feat/} sans \cmd{SPEC-NNN} est refusée avec le nommage attendu.
\end{itemize}
\astuce{Avec les hooks installés, tout commit sur \cmd{main} est refusé, y compris de documentation. Et \cmd{git commit --no-verify} contourne le hook : seule une CI côté serveur rendrait la règle absolue.}

\section{Les pull requests}
Le modèle \cmd{.github/pull_request_template.md} demande : la ou les specs concernées, l'ADR concerné, le lot ou la version, et les vérifications (spec validée, glossaire respecté, contrat à jour, un test par critère d'acceptation, fitness functions, branche et messages conformes).
